% Глава: 27-scripting-ide
% Статус: задел (кратко), будет дополнено.
\chapter{Скрипты оболочки и IDE}
\label{ch:27-scripting-ide}

\emph{Раздел будет дополнен в следующих версиях
пособия. Пока --- краткая вводка и ключевые темы.}

Скрипты и IDE мы используем на своей ВМ на \fpath{pglab.nsu.ru}:
собираем и запускаем проекты там, где они и живут.

\section{Скрипты оболочки}

Повторяющиеся команды логично собирать в скрипт
\fpath{run.sh}:
\begin{minted}[bgcolor=gray!6]{bash}
#!/usr/bin/env bash
# сборка и тестирование проекта
set -euo pipefail

make clean
make
make test
\end{minted}

\cmd{set -e} --- остановка при ошибке,
\cmd{set -u} --- ошибка на несуществующей
переменной, \cmd{set -o pipefail} --- ошибка в
пайпе, если любая команда пайпа упала.

\section{IDE для C и системной разработки}

\begin{itemize}
  \item \textbf{VS Code} + расширения (C/C++,
    clangd, CMake) --- популярен, настраиваем;
  \item \textbf{CLion} --- коммерческая IDE,
    мощный рефакторинг;
  \item \textbf{Neovim} + LSP --- для тех, кто
    хочет работать в терминале
    (\nameref{ch:09-vim}).
\end{itemize}

\section{Что будет в полном разделе}
\begin{itemize}
  \item Синтаксис bash: переменные, циклы,
    функции, условия;
  \item Аргументы командной строки (\cmd{\$1},
    \cmd{getopts});
  \item Обработка ошибок и логирование;
  \item Настройка IDE под C/C++: отладчик,
    компилятор, форматирование;
  \item LSP и автодополнение.
\end{itemize}

\section{Связь с остальным курсом}

Скрипты используют конвейеры
(\nameref{ch:07-text-pipelines}) и оболочки
(\nameref{ch:03-shell}); запуск в Docker
(\nameref{ch:19-docker-basics}); история изменений ---
Git (\nameref{ch:25-git}).

\section{Задачи для самопроверки}
\begin{enumerate}
  \item Зачем \cmd{set -euo pipefail} в скрипте?
  \item Как передать аргументы в bash-скрипт и
    как к ним обратиться?
  \item Чем IDE отличается от редактора текста?
\end{enumerate}
